% format.tex  --  macros, notation, styling. Swap the class without touching this.

\def\BibTeX{{\rm B\kern-.05em{\sc i\kern-.025em b}\kern-.08em
    T\kern-.1667em\lower.7ex\hbox{E}\kern-.125emX}}

% ---------------------------------------------------------------- draft switch
% \drafttrue  -> todo notes visible.   \draftfalse -> submission build.
\newif\ifdraft
\drafttrue
\ifdraft
  \newcommand{\todo}[1]{\textcolor{red}{[\textbf{TODO:} #1]}}
  \newcommand{\note}[1]{\textcolor{blue}{[#1]}}
\else
  \newcommand{\todo}[1]{}
  \newcommand{\note}[1]{}
\fi

% ------------------------------------------------------------ general notation
% Copied verbatim from scripts/gantry/orthogonal-by-construction/documentation/
% state-level-obc-derivation.tex so the two documents stay comparable.
\newcommand{\R}{\mathbb{R}}
\newcommand{\ran}{\operatorname{range}}
\newcommand{\rank}{\operatorname{rank}}
\newcommand{\col}{\operatorname{col}}
\newcommand{\norm}[1]{\left\lVert#1\right\rVert}
\newcommand{\code}[1]{\nolinkurl{#1}}

% ------------------------------------------------------------ paper notation
% Define the symbols for this paper here, once.
% Section II: gantry coordinates, baseline model, LFR.
\newcommand{\qact}{q_{\mathrm{act}}}     % actuator coordinates [X1 X2 Y]^T
\newcommand{\uact}{u_{\mathrm{act}}}     % actuator forces
\newcommand{\fnet}{f_{\mathrm{net}}}     % net generalized force on the inertia
\newcommand{\Dsch}{\Delta}               % scheduling block of the LFR
